problem:  
Let \((M^n,g,f)\) be a complete *asymptotically conical expanding gradient Ricci soliton*,  
\[
\mathrm{Ric}_g + \nabla^2 f = \lambda g, \qquad \lambda<0,
\]
asymptotic to the cone \(C(X)=(0,\infty)\times X\) with metric \(dr^2+r^2 h_0\), where \((X,h_0)\) is Einstein.  
Let \(L_\mathrm{AC}\) denote the linearization of the DeTurck‑gauged soliton operator at \((g,f)\), acting on symmetric \(2\)-tensors with weighted decay, and let \(L_{h_0}=\Delta_L+2(n-2)\) denote the Lichnerowicz operator on \(X\).

We know from elliptic AC theory that \(L_\mathrm{AC}\) is Fredholm in suitable weighted spaces away from indicial roots of \(L_{h_0}\).  

---

**Problem:**  
Beyond this structural Fredholmness, consider the *geometric meaning* of its kernel and image.

1. Does every \(\dot g\in\ker L_\mathrm{AC}\) correspond to a genuine infinitesimal deformation of the soliton that modifies the *asymptotic cone metric* at infinity?  
   Equivalently, is there a canonical surjective map
   \[
   \Phi : \ker L_\mathrm{AC} \longrightarrow \ker L_{h_0}
   \]
   describing the leading‐order asymptotic term of the deformation on the link?  If not, what analytic obstructions arise from the soliton potential \(f\)?

2. Are these infinitesimal cone deformations **integrable** at the nonlinear level—that is, do they generate genuine one‑parameter families of asymptotically conical expanders with perturbed link metrics \(h_0+t\,\dot h\)?  
   What analytic conditions, possibly involving weighted projections of quadratic error terms or the sign of the smallest Lichnerowicz eigenvalues, govern this integration problem?

3. Does the presence of a spectral gap in the positive spectrum of \(L_{h_0}\) guarantee *asymptotic stability* of the soliton under Ricci flow perturbations preserving the asymptotic cone type?

---

Why is it a "good" problem:  
This formulation pushes past known Fredholm properties to the **geometric consequences** of the analytic linearization: it asks how the kernel of a well‑understood elliptic operator governs real deformations of expanding Ricci solitons and their asymptotic cones, and when these infinitesimal modes can integrate into actual families of solitons.  

It thus bridges **spectral geometry on the conical link**, **analytic deformation theory**, and **Ricci flow dynamics**. The fundamental challenge lies in identifying and overcoming nonlinear obstructions introduced by the soliton potential, distinguishing between “true” geometric moduli of the asymptotic cone and pure gauge effects.  

The problem is simple to state yet profoundly deep—it aims to describe the local structure of the moduli space of asymptotically conical expanders and its connection to spectral data on the link, a question at the frontier of modern geometric analysis.